% Part: computability
% Chapter: computability-theory
% Section: prop-reduce

\documentclass[../../../include/open-logic-section]{subfiles}

\begin{document}

\olfileid{cmp}{thy}{ppr}
\olsection{归约的性质}

如果 $A$ 可归约到~$B$，我们就记作 $A \leq_m B$。这个记号暗示着，如果 $A \leq_m B$，那么 $A$ “不比 $B$ 难”，而 $B$ “至少和 $A$ 一样难”；并且 $\le_m$ 就像数之间通常的 $\le$ 一样，按复杂度对集合（或判定问题）排序。下面两个命题支持这一直观理解。第一个命题表明 $\le_m$ 具有传递性。

\begin{prop}
\ollabel{prop:trans-red}
如果 $A \leq_m B$ 且 $B \leq_m C$，那么 $A \leq_m C$。
\end{prop}

\begin{proof}
将 $A$ 到~$B$ 的归约与 $B$ 到 $C$ 的归约复合，就得到了 $A$ 到~$C$ 的归约。
\end{proof}

\begin{prob}
通过证明以下事实来证明 \olref{prop:trans-red}：如果 $f$ 和~$g$ 分别是 $A$ 到~$B$ 和 $B$ 到~$C$ 的多一归约，那么 $\comp{f}{g}$ 就是 $A$ 到~$C$ 的多一归约。
\end{prob}

\begin{prop}
\ollabel{prop:reduce}
设 $A$ 和 $B$ 是任意集合，并假设 $A \leq_m B$。
\begin{enumerate}
\item 如果 $B$ 是可计算枚举的，那么 $A$ 也是。
\item 如果 $B$ 是可计算的，那么 $A$ 也是。
\end{enumerate}
\end{prop}

\begin{proof}
设 $f$ 是从 $A$ 到~$B$ 的多一归约。对于第一个论断，只需验证：如果 $B$ 是某个偏函数~$g$ 的定义域，那么 $A$ 就是 $\comp{f}{g}$ 的定义域：
\begin{align*}
x \in A & \text{ 当且仅当 } f(x) \in B \\
& \text{ 当且仅当 }  g(f(x)) \fdefined.
\end{align*}

对于第二个论断，回忆如果 $B$~是可计算的，那么 $B$ 和~$\Complement{B}$ 都是可计算枚举的（\olref[cmp]{thm:ce-comp}）。不难验证 $f$~同样也是 $\Complement{A}$ 到 $\Complement{B}$ 的多一归约，因此，根据本证明的第一部分，$A$ 和~$\Complement{A}$ 都是可计算枚举的。所以由 \olref[cmp]{thm:ce-comp} 可知 $A$ 也是可计算的。
（或者，你可以验证 $\Char{A} = \comp{f}{\Char{B}}$；所以如果 $\Char{B}$ 是可计算的，那么 $\Char{A}$ 也是。）
\end{proof}

\begin{prob}
假设 $f$ 是 $A$ 到~$B$ 的多一归约。证明 $f$ 也是 $\Complement{A}$ 到 $\Complement{B}$ 的多一归约。
\end{prob}

\begin{prob}
证明：如果 $f\colon A \to B$ 是一个多一归约，那么 $\Char{A}
= \comp{f}{\Char{B}}$。
\end{prob}

\begin{digress}
一种更一般的归约概念，称为\emph{Turing归约}，在其他语境中很有用，特别是用于证明不可判定性结果。注意，根据\olref[cmp]{cor:comp-k}，$K_0$ 的补集不能归约到~$K_0$，因为它不是可计算枚举的。但是，直观上，如果你知道关于 $K_0$ 的问题的答案，你也就知道了关于其补集的问题的答案。如果可以使用一个能询问关于~$B$ 的问题的可计算程序来确定~$A$ 中问题的答案，则称集合 $A$ 可\emph{Turing归约}到~$B$。这比多一归约更为宽松，因为在多一归约中，(1)~只允许询问一个关于 $B$ 的问题，并且 (2)~关于 $B$ 的“是”答案必须转化为关于 $A$ 的问题的“是”答案，对“否”亦然。仍然成立的是，如果 $A$~可Turing归约到~$B$ 并且 $B$~是可计算的，那么 $A$~也是可计算的（不过，正如我们所见，对于可计算枚举性，类似的结论并不成立）。

你应该思考我们讨论过的各种归约概念，并理解它们之间的区别。然而，在本章中，我们只处理多一归约。顺便提一下，上一段讨论的这两种归约在计算复杂性中都有对应物，只是附加了Turing机必须在多项式时间内运行的要求：多一归约的复杂性版本称为\emph{Karp（卡普）归约}，而Turing归约的复杂性版本称为\emph{Cook（库克）归约}。
\end{digress}

\end{document}